% Informe del entregable — arquitectura medallón.
% Compilar:  latexmk -xelatex informe.tex     (o: xelatex informe.tex, dos veces)
%
% ── Datos de la portada: editá SOLO estas cuatro líneas ─────────────────────
\newcommand{\Autora}{Erika Jimenez}
\newcommand{\Materia}{Práctica Módulo 1}
\newcommand{\Institucion}{Diplomado de AI · Augmented Humans}
\newcommand{\Fecha}{1 de agosto de 2026}
% ───────────────────────────────────────────────────────────────────────────

\documentclass[11pt,a4paper]{article}

\usepackage{fontspec}
\setmainfont{Latin Modern Roman}
\setmonofont{DejaVu Sans Mono}[Scale=0.82]
\usepackage[spanish,es-noquoting]{babel}
\usepackage[margin=2.4cm]{geometry}
\usepackage{booktabs}
\usepackage{longtable}
\usepackage{array}
\usepackage{xcolor}
\usepackage{fvextra}
\usepackage[most]{tcolorbox}
\usepackage{tikz}
\usetikzlibrary{arrows.meta,positioning,shapes.geometric}
\usepackage{listings}
\usepackage{caption}
\usepackage{hyperref}

\definecolor{bronce}{HTML}{A8703E}
\definecolor{plata}{HTML}{6E7B8B}
\definecolor{oro}{HTML}{B8912F}
\definecolor{tinta}{HTML}{1F2933}
\definecolor{suave}{HTML}{F4F5F7}
\definecolor{acento}{HTML}{2C5F8A}

\hypersetup{colorlinks=true, linkcolor=acento, urlcolor=acento, citecolor=acento}

\lstdefinestyle{codigo}{
  basicstyle=\ttfamily\footnotesize,
  backgroundcolor=\color{suave},
  keywordstyle=\color{acento}\bfseries,
  commentstyle=\color{plata}\itshape,
  stringstyle=\color{bronce},
  breaklines=true,
  showstringspaces=false,
  frame=leftline,
  framerule=1.6pt,
  rulecolor=\color{acento!40},
  xleftmargin=10pt,
  framexleftmargin=8pt,
  aboveskip=10pt, belowskip=10pt,
  captionpos=b,
}
\lstset{style=codigo}
\lstdefinelanguage{sqlite}{
  morekeywords={INSERT,INTO,SELECT,FROM,WHERE,ON,CONFLICT,DO,UPDATE,SET,NOTHING,
    CREATE,OR,REPLACE,TABLE,AS,QUALIFY,PARTITION,BY,ORDER,GROUP,HAVING,COUNT,JOIN,
    USING,PRIMARY,KEY,DISTINCT,CASE,WHEN,THEN,ELSE,END,IS,NOT,excluded},
  sensitive=false, morecomment=[l]{--},
}

\newtcolorbox{caja}[1][]{colback=suave, colframe=acento!55, boxrule=0.7pt,
  arc=2pt, left=8pt, right=8pt, top=6pt, bottom=6pt, #1}

\captionsetup{font=small, labelfont={bf,color=acento}}
\setlength{\parskip}{0.45em}
\setlength{\parindent}{0pt}

\title{\vspace{-1.5cm}\textbf{Arquitectura medallón para un corpus de noticias
de mercado}\\[0.3em]\large Bronze crudo, contrato de datos, carga idempotente
y búsqueda semántica vectorial}
\author{\Autora \\ \small \Materia\ ---\ \Institucion}
\date{\Fecha}

\begin{document}
\maketitle

\begin{caja}
\textbf{Resumen.} Se implementó una arquitectura medallón completa sobre
noticias de mercado cripto, ingeridas de tres feeds RSS públicos sin
autenticación. La capa \emph{Bronze} conserva dos lotes crudos íntegros con
timestamp y huella SHA-256; \emph{Silver} valida cada registro contra un
contrato Pydantic v2 y desvía los inválidos a una tabla de cuarentena con el
motivo del rechazo; la carga usa \emph{staging} más \texttt{UPSERT} por clave
natural, de modo que reprocesar un lote arroja \texttt{filas\_nuevas = 0}; y
\emph{Gold} construye un índice vectorial FAISS de 384 dimensiones que expone
una función de búsqueda semántica. Todo corre en contenedor vía
\texttt{docker-compose} y fue ejecutado de punta a punta cuatro veces. Los cinco
criterios de aceptación se cumplen con evidencia reproducible.
\end{caja}

\section{Objetivo y dominio}

El requisito pedía una arquitectura medallón sobre un tema de interés propio,
alimentada por una API o feed público real y sin autenticación. El dominio
elegido son las \textbf{noticias de mercado de criptoactivos}, por continuidad
con un proyecto propio de investigación cuantitativa: el corpus de noticias es
texto libre, heterogéneo y sucio, que es justamente donde un contrato de datos y
una búsqueda semántica aportan valor real y no decorativo.

Las fuentes son tres feeds RSS abiertos ---CoinDesk, Cointelegraph y
Decrypt---, sin API key, sin token y sin cuenta. El entregable es una base de
código autocontenida (1\,176 líneas de Python, 15 pruebas automatizadas) que no
depende de ningún servicio de pago.

\section{Arquitectura}

\begin{center}
\resizebox{\textwidth}{!}{%
\begin{tikzpicture}[
  font=\small,
  capa/.style={rectangle, rounded corners=3pt, draw, thick, minimum width=3.5cm,
               minimum height=1.5cm, align=center, text=white},
  nota/.style={align=center, font=\scriptsize\itshape, text=tinta!70},
  flecha/.style={-{Stealth[length=3mm]}, thick, draw=tinta!60}
]
\node[capa, fill=bronce] (b) {\textbf{BRONZE}\\[2pt]\scriptsize payload crudo\\\scriptsize + SHA-256};
\node[capa, fill=plata, right=2.1cm of b] (s) {\textbf{SILVER}\\[2pt]\scriptsize contrato Pydantic\\\scriptsize + UPSERT};
\node[capa, fill=oro, right=2.1cm of s] (g) {\textbf{GOLD}\\[2pt]\scriptsize embeddings\\\scriptsize + índice FAISS};

\node[left=1.5cm of b, align=center, font=\scriptsize] (rss) {3 feeds RSS\\públicos\\(sin auth)};
\node[below=1.05cm of s, align=center, font=\scriptsize, draw, dashed, rounded corners=2pt,
      inner sep=4pt, text=tinta] (q) {cuarentena\\\texttt{silver\_rejects}\\(campo + motivo)};
\node[below=1.05cm of g, align=center, font=\scriptsize, draw, rounded corners=2pt,
      inner sep=4pt, text=tinta] (api) {\texttt{search(query, k)}\\coseno top-$k$};

\draw[flecha] (rss) -- (b);
\draw[flecha] (b) -- node[above, nota] {parse} (s);
\draw[flecha] (s) -- node[above, nota] {embed} (g);
\draw[flecha] (g) -- (api);
\draw[flecha, draw=tinta!45] (s) -- node[right, nota, xshift=2pt] {inválidos} (q);

\node[below=0.35cm of b, nota] {inmutable};
\node[below=2.7cm of s, nota, yshift=0.2cm] {DuckDB · \texttt{data/medallion.duckdb}};
\end{tikzpicture}}
\end{center}

El almacenamiento es \textbf{DuckDB} en archivo, montado como volumen del
contenedor: el estado sobrevive entre corridas, condición necesaria para poder
demostrar idempotencia. Las tablas son seis:

\begin{center}
\small
\begin{tabular}{@{}llp{7.6cm}@{}}
\toprule
\textbf{Capa} & \textbf{Tabla} & \textbf{Rol} \\
\midrule
Bronze & \texttt{bronze\_batches} & payload crudo por (lote, fuente) + SHA-256 \\
Silver & \texttt{stg\_news} & \emph{staging}, se recrea en cada corrida \\
Silver & \texttt{silver\_news} & destino limpio, PK = clave natural \\
Silver & \texttt{silver\_rejects} & cuarentena con campo, tipo y motivo \\
Gold & \texttt{gold\_news\_embeddings} & mapa \texttt{news\_id} $\rightarrow$ posición en FAISS \\
--- & \texttt{load\_audit} & bitácora por corrida (evidencia de idempotencia) \\
\bottomrule
\end{tabular}
\end{center}

\section{Bronze: ingesta cruda}

La regla de la capa es que el payload se guarde \textbf{tal como llegó}: sin
parsear, sin limpiar HTML y sin recortar. Se persiste por partida doble ---
archivo en \texttt{data/medallion/bronze/<batch\_id>/<fuente>.xml} y renglón en
\texttt{bronze\_batches}--- con su huella SHA-256, de modo que cualquiera puede
verificar después que el contenido no fue alterado. El \texttt{batch\_id} es el
timestamp UTC de la corrida, y la escritura usa \texttt{ON CONFLICT DO NOTHING}:
un lote ya escrito nunca se sobrescribe.

\begin{caja}
\textbf{Por qué no limpiar aquí.} La tentación es normalizar el HTML en la
ingesta. Si se hace, se pierde la capacidad de re-parsear el histórico con
reglas nuevas cuando el contrato cambie ---y con ella, el sentido de tener una
capa Bronze. La limpieza vive en Silver.
\end{caja}

\begin{table}[h]
\centering\small
\caption{Lotes crudos almacenados (salida de \texttt{bronze\_batches}).}
\begin{tabular}{@{}lrlrr@{}}
\toprule
\textbf{batch\_id} & \textbf{fuentes} & \textbf{fetched\_at} & \textbf{bytes crudos} & \textbf{SHA-256 distintos} \\
\midrule
\texttt{20260801T001833Z} & 3 & 2026-08-01 00:18:33 & 116\,368 & 3 \\
\texttt{20260801T001859Z} & 3 & 2026-08-01 00:18:59 & 116\,368 & 3 \\
\bottomrule
\end{tabular}
\end{table}

\section{Silver: contrato de datos y cuarentena}

Cada registro se valida contra el modelo \texttt{NewsItem} (Pydantic v2) antes
de tocar la tabla destino. El contrato es explícito y estricto: declara
\texttt{extra="forbid"}, de modo que un campo inesperado en el feed es un
rechazo y no un silencio.

\begin{lstlisting}[language=Python, caption={Contrato de datos (extracto de \texttt{contracts.py}).}]
class NewsItem(BaseModel):
    model_config = ConfigDict(extra="forbid", str_strip_whitespace=True)

    news_id:      Annotated[str, Field(min_length=32, max_length=32)]
    source:       Literal["coindesk", "cointelegraph", "decrypt"]
    title:        Annotated[str, Field(min_length=10,  max_length=300)]
    url:          Annotated[str, Field(min_length=12,  max_length=2000)]
    body:         Annotated[str, Field(min_length=40,  max_length=8000)]
    published_at: datetime
    symbols:      list[str] = Field(default_factory=list)

    @field_validator("published_at")
    @classmethod
    def _fecha_plausible(cls, v: datetime) -> datetime:
        # ni del futuro ni de hace mas de tres anos
        ...

    @model_validator(mode="after")
    def _id_consistente(self) -> "NewsItem":
        if self.news_id != natural_key(self.source, self.url):
            raise ValueError("news_id no corresponde a sha256(source|url)")
        return self
\end{lstlisting}

Lo que no pasa el contrato no se descarta: se escribe en
\texttt{silver\_rejects} junto con el \textbf{campo}, el \textbf{tipo de error},
el \textbf{mensaje} y el \textbf{registro crudo completo} en JSON, de manera que
el rechazo sea auditable y reprocesable. Cuando fallan varios campos a la vez se
reporta la \emph{causa raíz}: si la URL falta, \texttt{news\_id} también falla
por ser derivado, y el motivo relevante es el primero.

\begin{table}[h]
\centering\small
\caption{Validación por lote. Los cuatro rechazos (2 por lote) comparten motivo:
campo \texttt{body}, tipo \texttt{string\_too\_short}, mensaje \emph{«String should
have at least 40 characters»} ---notas cuyo resumen RSS venía vacío o de una
sola línea---.}
\begin{tabular}{@{}lrrr@{}}
\toprule
\textbf{batch\_id} & \textbf{leídas} & \textbf{válidas} & \textbf{rechazadas} \\
\midrule
\texttt{20260801T001833Z} & 93 & 91 & 2 \\
\texttt{20260801T001859Z} & 93 & 91 & 2 \\
\bottomrule
\end{tabular}
\end{table}

\section{Carga idempotente}

La clave natural del dominio es \texttt{news\_id = sha256(source | url)[:32]}, y
es \texttt{PRIMARY KEY} de \texttt{silver\_news}. Se eligió la URL y no el
título ---que cambia con las ediciones del editor--- ni el GUID del feed ---que
cada fuente formatea distinto---.

El flujo es \emph{staging} y luego MERGE. La tabla \texttt{stg\_news} se recrea
en cada corrida y se deduplica internamente, porque un mismo feed puede repetir
una nota dentro del mismo lote; sin ese paso el \texttt{UPSERT} intentaría
actualizar dos veces la misma fila.

\begin{lstlisting}[language=sqlite, caption={Deduplicación en staging y MERGE por clave natural (\texttt{silver.py}).}]
CREATE OR REPLACE TABLE stg_news AS
SELECT * FROM stg_news
QUALIFY row_number() OVER (PARTITION BY news_id ORDER BY published_at DESC) = 1;

INSERT INTO silver_news (news_id, source, title, url, body, published_at,
                         symbols, content_hash, first_batch_id, last_batch_id,
                         inserted_at, updated_at)
SELECT news_id, source, title, url, body, published_at, symbols,
       content_hash, batch_id, batch_id, ?, ?
FROM stg_news
ON CONFLICT (news_id) DO UPDATE SET
    title         = excluded.title,
    body          = excluded.body,
    published_at  = excluded.published_at,
    content_hash  = excluded.content_hash,
    last_batch_id = excluded.last_batch_id,
    updated_at    = CASE
                        WHEN silver_news.content_hash IS DISTINCT FROM excluded.content_hash
                        THEN excluded.updated_at
                        ELSE silver_news.updated_at
                    END;
\end{lstlisting}

El \texttt{CASE} sobre \texttt{content\_hash} ---resumen de título, cuerpo, fecha
y símbolos--- distingue dos situaciones que suelen confundirse: \emph{«ya había
visto esta nota»} y \emph{«esta nota cambió»}. Un reproceso puro deja
\texttt{updated\_at} intacto; una corrección editorial en la fuente sí queda
registrada. Cada corrida escribe su renglón en \texttt{load\_audit}.

\begin{table}[h]
\centering\small
\caption{Bitácora \texttt{load\_audit}: cuatro ejecuciones de punta a punta.}
\begin{tabular}{@{}llrrr@{}}
\toprule
\textbf{run\_id} & \textbf{batch\_id} & \textbf{válidas} & \textbf{filas\_nuevas} & \textbf{filas\_actualizadas} \\
\midrule
\texttt{run-20260801T001834Z} & \texttt{...T001833Z} & 91 & 91 & 0 \\
\texttt{run-20260801T001900Z} & \texttt{...T001859Z} & 91 & \textbf{0} & 0 \\
\texttt{run-20260801T001939Z} & \texttt{...T001833Z} & 91 & \textbf{0} & 0 \\
\texttt{run-20260801T001953Z} & \texttt{...T001859Z} & 91 & \textbf{0} & 0 \\
\bottomrule
\end{tabular}
\end{table}

Las corridas 3 y 4 son el caso estricto que pide el criterio: reprocesar el
\emph{mismo} \texttt{batch\_id} ya cargado. Ambas dan \texttt{filas\_nuevas = 0}.
La corrida 2 es un caso adicional e interesante: bajó un lote \emph{nuevo}
(timestamp y SHA-256 propios) veintiséis segundos después del primero; como los
feeds no habían publicado nada en ese lapso, el MERGE reconoció las 91 notas por
clave natural y tampoco insertó nada. La deduplicación es por contenido, no por
lote.

\section{Ausencia de duplicados}

\begin{lstlisting}[language=sqlite, caption={Consulta de verificación y su resultado.}]
SELECT news_id, count(*) FROM silver_news GROUP BY news_id HAVING count(*) > 1;
-- (0 filas)
-- filas en silver_news = 91  ·  news_id distintos = 91
\end{lstlisting}

La cuarentena también es idempotente: su llave primaria es
\texttt{(batch\_id, source, record\_ord)}, así que reprocesar un lote no infla la
tabla de rechazos ---un detalle que se escapa con facilidad y que está cubierto
por una prueba automatizada.

\section{Gold: índice vectorial y búsqueda semántica}

El campo relevante es el texto de la nota: se embebe \texttt{título. cuerpo} con
\texttt{sentence-transformers/all-MiniLM-L6-v2} (384 dimensiones, CPU, local) y
se indexa en un \texttt{faiss.IndexFlatIP} sobre vectores normalizados en L2, de
modo que el producto interno \textbf{es} la similitud coseno. El índice se
persiste en disco y la construcción es incremental: sólo se embeben los
\texttt{news\_id} que aún no están en \texttt{gold\_news\_embeddings}, cuya
columna \texttt{vec\_ord} guarda la posición en el índice.

Se eligió el índice exhaustivo y no uno aproximado (IVF, HNSW) porque a esta
escala ---decenas o centenas de documentos--- la búsqueda exacta es instantánea
y un índice aproximado sólo agregaría hiperparámetros que calibrar sin ganancia
medible.

\begin{lstlisting}[language=Python, caption={Función de búsqueda semántica expuesta (\texttt{gold.py}).}]
def search(query: str, k: int = 5) -> list[dict]:
    """Devuelve las k notas mas cercanas al texto libre."""
    idx = _load_index()
    qv = embed([query])                 # normalizado en L2
    scores, ords = idx.search(qv, min(k, idx.ntotal))
    # vec_ord -> JOIN con silver_news -> filas con score coseno
\end{lstlisting}

\begin{table}[h]
\centering\small
\caption{Estado del índice tras las cuatro corridas.}
\begin{tabular}{@{}rrlr@{}}
\toprule
\textbf{vectores} & \textbf{vec\_ord distintos} & \textbf{modelo} & \textbf{dim} \\
\midrule
91 & 91 & \texttt{all-MiniLM-L6-v2} & 384 \\
\bottomrule
\end{tabular}
\end{table}

\subsection*{Consultas demostradas}

\begin{table}[h]
\centering\small
\caption{Tres consultas en lenguaje natural y sus tres primeros resultados.}
\begin{tabular}{@{}rp{11.5cm}@{}}
\toprule
\textbf{score} & \textbf{titular recuperado} \\
\midrule
\multicolumn{2}{@{}l@{}}{\textit{«lawsuit and regulatory crackdown by financial authorities»}} \\
0.4627 & New York \textbf{sues} Kalshi over alleged illegal gambling operation \\
0.4091 & New York sues Kalshi, alleges it offers a gambling platform «plain and...» \\
0.3604 & The good and the bad of perps, according to crypto traders \\
\midrule
\multicolumn{2}{@{}l@{}}{\textit{«institutional money flowing into spot ETFs»}} \\
0.5666 & Bitcoin ETFs post \$233M \textbf{inflows}, pushing week back into the green \\
0.5391 & Bitcoin, Ethereum Wobble as Fed Holds Rates Steady \\
0.4162 & Aviva Investors launches tokenized fund after Central Bank of Ireland... \\
\midrule
\multicolumn{2}{@{}l@{}}{\textit{«sharp market selloff and price volatility»}} \\
0.4124 & Bitcoin's calm is back and so is the setup for a \textbf{volatility explosion} \\
0.3029 & Bitcoin, ether fall, equities rally with broader crypto market on track... \\
0.3014 & Bitcoin holds monthly gain, faces «choppy» August as «forced-selling»... \\
\bottomrule
\end{tabular}
\end{table}

\begin{caja}
\textbf{Por qué es búsqueda semántica y no textual.} La consulta \emph{«lawsuit
and regulatory crackdown by financial authorities»} recupera \emph{«New York
sues Kalshi...»}: ninguna de las palabras de la consulta aparece en ese titular
---un \texttt{LIKE '\%lawsuit\%'} habría devuelto cero filas---. Lo mismo ocurre
con \emph{«institutional money flowing into spot ETFs»}, que encuentra
\emph{«inflows»}, y con \emph{«sharp market selloff»}, que encuentra
\emph{«volatility explosion»}. La recuperación opera sobre el significado del
texto, no sobre sus caracteres.
\end{caja}

\section{Contenedorización y reproducción}

La imagen pesa 459\,MB: se instala explícitamente la build \emph{CPU} de PyTorch
desde su índice, porque la build por defecto arrastra CUDA (unos 2,5\,GB) que
aquí no se usa. El contenedor corre como UID 1000 para que lo escrito en el
volumen pertenezca al usuario del anfitrión y no a \texttt{root}.

\begin{lstlisting}[language=bash, caption={Reproducción completa desde cero.}]
cd medallion
docker compose build

docker compose run --rm medallion                    # corrida 1: lote nuevo
docker compose run --rm medallion                    # corrida 2: segundo lote
docker compose run --rm medallion python -m medallion run --batch-id <ID>
docker compose run --rm medallion python -m medallion search "texto libre"
\end{lstlisting}

El directorio \texttt{data/} se monta como volumen: bronze crudo, base DuckDB e
índice FAISS sobreviven al contenedor, que es precisamente lo que permite
demostrar idempotencia entre ejecuciones separadas.

\section{Pruebas automatizadas}

Quince pruebas corren sin red y sin descargar el modelo ---los embeddings se
sustituyen por un hash determinista y el lote de Bronze se siembra con un RSS
sintético---:

\begin{itemize}\setlength\itemsep{0.1em}
  \item aceptación del contrato y rechazo campo por campo (URL inválida, cuerpo
        corto, título corto, fuente desconocida, fecha futura, símbolo fuera de
        la whitelist, campo extra, \texttt{news\_id} inconsistente);
  \item cuarentena que registra el motivo correcto;
  \item triple reproceso del mismo lote con \texttt{filas\_nuevas = 0};
  \item dos lotes de contenido idéntico sin generar duplicados;
  \item cuarentena que no se infla al reprocesar;
  \item índice FAISS incremental y orden decreciente de \emph{score} en la
        búsqueda.
\end{itemize}

\begin{lstlisting}[language=bash]
$ pytest medallion/tests -o addopts=""
15 passed in 1.21s
\end{lstlisting}

\section{Cumplimiento de los criterios de aceptación}

\begin{center}
\small
\begin{tabular}{@{}p{2.3cm}p{5.6cm}p{7.4cm}@{}}
\toprule
\textbf{Criterio} & \textbf{Mínimo exigido} & \textbf{Obtenido} \\
\midrule
Bronze & 2+ lotes crudos, timestamp, sin transformar &
2 lotes de 116\,368 bytes, 3 payloads con SHA-256 cada uno, guardados byte a
byte en disco y en base \\
\addlinespace
Contrato & validación explícita, cuarentena con motivo &
\texttt{NewsItem} (Pydantic v2, \texttt{extra="forbid"}); 93 leídas $\rightarrow$
91 válidas y 2 en cuarentena por lote, con campo, tipo, mensaje y registro crudo \\
\addlinespace
Idempotencia & reproceso da \texttt{filas\_nuevas = 0} &
staging + \texttt{ON CONFLICT DO UPDATE}; 3 de 4 corridas con
\texttt{filas\_nuevas = 0}, incluidos dos reprocesos del mismo \texttt{batch\_id} \\
\addlinespace
Duplicados & \texttt{COUNT(*) > 1} por clave natural = 0 filas &
0 filas; 91 registros con 91 \texttt{news\_id} distintos \\
\addlinespace
Gold & índice vectorial funcional, 1+ consulta semántica &
FAISS \texttt{IndexFlatIP} de 384-d con 91 vectores; 3 consultas demostradas \\
\bottomrule
\end{tabular}
\end{center}

\section{Decisiones de diseño y limitaciones}

Tres decisiones merecen mención explícita. Primera, la clave natural es la URL y
no el título ni el GUID, por las razones ya expuestas. Segunda, Bronze es
inmutable por diseño (\texttt{ON CONFLICT DO NOTHING}), lo que permite re-parsear
el histórico si el contrato evoluciona. Tercera, el \texttt{content\_hash}
separa la re-observación del cambio real, que es lo que hace que la métrica
\texttt{filas\_actualizadas} signifique algo.

Las limitaciones se enuncian sin adorno:

\begin{itemize}\setlength\itemsep{0.1em}
  \item el modelo \texttt{all-MiniLM-L6-v2} es \textbf{monolingüe en inglés},
        igual que el corpus; una consulta en español recupera resultados
        notablemente peores, y una versión bilingüe exigiría cambiar a
        \texttt{paraphrase-multilingual-MiniLM-L12-v2};
  \item los feeds RSS entregan sólo el resumen y no la nota completa, de modo
        que los embeddings describen el \emph{abstract}, no el cuerpo entero;
  \item la inferencia de símbolos es por palabras clave y no por reconocimiento
        de entidades: alcanza para etiquetar, no para desambiguar.
\end{itemize}

\appendix
\section{Anexo: salida literal de la ejecución}

Lo que sigue es la salida sin editar de
\texttt{docker compose run -{}-rm medallion python -m medallion evidence},
ejecutada tras las cuatro corridas.

{\footnotesize
\VerbatimInput[breaklines=true, breakanywhere=true, fontsize=\scriptsize]{salida-e2e.txt}
}

\end{document}
